\subsection{Chain-of-Thought for Spatial Reasoning}
\label{subsec:p3_s7_3_3}
\label{ssec:cot-reasoning}

Spatial reasoning poses unique challenges for large language models (LLMs), as it requires the manipulation of geometric relations, perspective-taking, and dynamic transformations: skills that are not purely symbolic but grounded in embodied cognition.
Chain-of-Thought (CoT) prompting~\cite{wei2022chain}, originally developed as a mechanism to elicit intermediate reasoning steps in arithmetic and commonsense tasks, has therefore been imported into spatial reasoning not as a generic recipe for longer answers, but as a way to externalize the intermediate spatial commitments that a model must keep consistent: which object is left of which, which frame of reference is active, which action changes the scene, and which hypothesized 3D configuration explains the observation.
The evidence is mixed.
In some settings, explicit rationales help decompose a spatial query into auditable subgoals; in others, they induce a fluent but ungrounded verbal trace that makes spatial performance worse~\cite{kancheti2026cotdegrades}.

\begin{figure}[t]
    \centering
    \includegraphics[width=\linewidth]{figures/part3/section7/chain-of-thought_spatial_reasoning.pdf}
    \caption{Spatial reasoning tasks (top) and intermediate-trace structures (bottom). Top: (a) textual deduction in a north-up frame, (b) visual counting, where the boxed red cube is the only one left of the green sphere, and (c) embodied planning revised by new observations. Bottom: CoT, ToT, and GoT differ in trace topology; PoT and CoA ground the trace in program execution and environment interaction. $\tau=(t_1,\ldots,t_L)$ is the intermediate trace and $\hat y=g(c,q,\tau)$ the final answer. Scenes and trace diagrams are illustrative, not recorded model computation.}
    \label{fig:p3_s7_3_3_cot_formulation}
\end{figure}


\subsubsection{Task Definition}
The task is to answer a spatial query $q$ about a scene or environment, given context $c$ that may be purely textual, visual, 3D, temporal, or embodied.
A CoT-style reasoner augments the answer with an intermediate trace
\begin{equation}
    \tau = (t_1,\ldots,t_L), \qquad \hat{y}=g(c,q,\tau),
\end{equation}
where each $t_\ell$ records a partial spatial commitment: a relation, coordinate estimate, subgoal, viewpoint change, object binding, or predicted action effect.
The trace is useful only insofar as it improves the grounded answer $\hat{y}$ or makes the failure mode inspectable; many apparent ``reasoning steps'' are not grounded in metric or perceptual evidence and merely rationalize a final guess, so the central question is not whether a model can produce more text, but whether the trace helps preserve spatial invariants across transformations, viewpoints, time steps, and actions.

\subsubsection{Methodology}
CoT, Tree-of-Thought (ToT)~\cite{yao2023tot}, and Graph-of-Thought (GoT)~\cite{besta2024got} differ mainly in the topology of the intermediate trace: CoT is a single path through partial spatial commitments; ToT branches over alternative hypotheses or subplans and uses a value function to prune; GoT allows shared sub-derivations, such as reusing a computed room layout or object graph across several queried relations.
All three search over partial reasoning states with a proposal mechanism, an evaluator $V$, and a halting rule $H$,
\begin{equation}
    \tau^\star=\arg\max_{\tau} V(\tau)\quad\text{s.t.}\quad t_{\ell+1}\in\mathrm{Propose}(t_{\le \ell}),\ \ H(\tau)=1,
\end{equation}
and differ in whether the search structure is a path, a tree of alternatives, or a graph of reusable sub-results, and in whether the trace stays verbal or becomes executable through programs (Program-of-Thought, PoT~\cite{chen2022pot}) or environment-changing actions (Chain-of-Action, CoA~\cite{pan2024coa}).
The use cases below organize how these traces are deployed.

\paragraph{Text-Only Spatial Deduction.}
In linguistic spatial puzzles, route descriptions, and relative-position QA, the model must maintain a latent arrangement without visual feedback.
CoT can expose intermediate bindings (\eg, ``A is north of B; C is west of A'') and reduce one-shot relational confusion~\cite{hu2023chain}, while ToT-style branching can preserve competing layouts until enough constraints have been applied.
The persistent failure mode is coordinate drift: even with explicit steps, LLMs lose track of the origin, orientation, or accumulated displacement in textual mazes and grid-navigation tasks~\cite{zhang2025mitigating}, so text-only CoT is most reliable when the trace is tied to a compact symbolic artifact (a table, coordinate sketch, or adjacency list) rather than to unconstrained prose.

\paragraph{Multimodal Scene Reasoning.}
In image- or video-grounded QA, a CoT trace must bind language to perceptual evidence: which object instance is referred to, what relation is visible, which is occluded, and which viewpoint assumption is made.
Free-form rationales can verbalize a relation that the image does not support, so recent work moves toward grounded intermediate artifacts (structured layouts, object bindings, diagrams, or executable module calls): SCoT-style planners use layout-coordinate traces to guide diffusion models on the generation side~\cite{chen2026spatialcotgeneration}, and program-mediated systems call perception modules rather than describe them~\cite{hsu2025makes}, the paradigm of §\ref{subsec:p3_s7_3_4}.
Here ToT and GoT serve less as generic search algorithms than as hypothesis-management devices that let a model compare object correspondences, camera frames, or depth-order explanations before committing to an answer.

\paragraph{Embodied Planning and VLA Reasoning.}
In navigation, manipulation, and vision-language-action tasks, spatial reasoning unfolds over time, and the trace must connect perceptual state, subgoal choice, action feasibility, and feedback from the environment.
CoT-style subgoals support long-horizon planning in navigation~\cite{lin2025navcot,liu2025navr1}, embodied instruction following~\cite{lu2023thinkbot,mu2023embodiedgpt}, and VLA models that expose intermediate visual or action commitments: CoT-VLA pairs reasoning with predicted visual states~\cite{zhao2025cotvla}, ACoT-VLA formulates the trace in action space through coarse action intents and reference trajectories~\cite{zhong2026acotvla}, and scene graphs stabilize the trace by giving the model a persistent substrate for object identities and relations~\cite{zhang2025embodiedvsr}.

\subsubsection{Benchmarks}
Evaluation mirrors the three use cases.
Spatial puzzles, relative-position QA, and textual maze navigation are evaluated with prompt-based CoT on the suites of~\cite{hu2023chain,zhang2025mitigating}; visual question answering and scene-grounded compositional queries are evaluated under CoT with program execution (§\ref{subsec:p3_s7_3_4}), where the rationale's correctness is mediated by program output rather than by the rationale string; and long-horizon planning with sensory feedback is evaluated through CoT-augmented navigation (NavCoT~\cite{lin2025navcot}, NavR1~\cite{liu2025navr1}), CoT-augmented instruction following (ThinkBot~\cite{lu2023thinkbot}, EmbodiedGPT~\cite{mu2023embodiedgpt}), visual- and action-trace supervision for VLA models~\cite{zhao2025cotvla,zhong2026acotvla}, and scene-graph-augmented embodied reasoning~\cite{zhang2025embodiedvsr}.


\subsubsection{Challenges and Trends}
The transfer of CoT to spatial intelligence is uneven.
\textbf{Coordinate hallucination} persists under simple textual CoT: even with ToT-style beam expansion, prompt-based planners mistrack absolute positions across multi-step chains in maze and grid-navigation problems~\cite{zhang2025mitigating}.
\textbf{Ungrounded supervision} can make things worse: textual CoT applied without a structured or rationale-grounded training signal degrades spatial VLM performance, because the model is forced to verbalize a reasoning trace it does not possess~\cite{ji2025enhancing,kancheti2026cotdegrades}, especially when the prompt rewards plausible explanation length over geometric correctness.
\textbf{Search cost} scales super-linearly in solution depth for ToT and GoT (and in the branching factor for ToT), which limits real-time or on-device deployment.
\textbf{Trace--action mismatch} arises in embodied settings when a VLA model produces a coherent verbal rationale while its action head follows a different latent policy; action-aware CoT aligns rationales with action tokens and visual predictions~\cite{zhong2026acotvla}, but reliable attribution between the visible trace and the executed behavior remains unresolved.

\emph{Grounded rationale supervision} pairs CoT with explicit spatial signals (scene graphs~\cite{zhang2025embodiedvsr}, depth maps, 3D coordinate frames, action traces, or layout-coordinate traces for spatial generation~\cite{chen2026spatialcotgeneration}) so the rationale is verifiable rather than free-form.
A complementary line makes the trace \emph{execution-grounded}: PoT-style traces compile and run against a geometric engine or vision--language tool, with execution outputs serving as a truth signal, the natural bridge from CoT to code / DSL generation in §\ref{subsec:p3_s7_3_4}; for embodied agents the parallel bridge runs from CoT to CoA and the agentic orchestration loops of §\ref{subsec:p3_s7_3_7}.
